% ═══════════════════════════════════════════════════════════════
% MT-05d-ML: MUSTERLÖSUNG WIEDERHOLUNGSTEST MI TIEMPO LIBRE
% Spanisch Klasse 7 - Sequenz 5d (Abschlusstest)
% ═══════════════════════════════════════════════════════════════
% Punkte: 30 BE
% ═══════════════════════════════════════════════════════════════

\documentclass[11pt, a4paper]{article}

\newif\ifswmodus
\swmodusfalse

\usepackage[utf8]{inputenc}
\usepackage[T1]{fontenc}
\usepackage[ngerman]{babel}
\usepackage{helvet}
\renewcommand{\familydefault}{\sfdefault}

\usepackage[margin=2cm]{geometry}
\usepackage{enumitem}
\usepackage{tabularx}
\usepackage{booktabs}

\usepackage{xcolor}
\usepackage{tcolorbox}
\tcbuselibrary{skins}

\usepackage{fancyhdr}
\usepackage{lastpage}
\usepackage{needspace}

% FARBEN
\definecolor{spanisch-color}{HTML}{c41e3a}
\definecolor{grau-color}{HTML}{888888}
\definecolor{hellgrau-color}{HTML}{f5f5f5}
\definecolor{loesung-color}{HTML}{c62828}
\definecolor{er-color}{HTML}{2c5aa0}
\definecolor{ir-color}{HTML}{2a7a4b}

\definecolor{sw-dark}{HTML}{000000}
\definecolor{sw-gray}{HTML}{666666}
\definecolor{sw-white}{HTML}{ffffff}

\ifswmodus
    \colorlet{spanisch}{sw-dark}
    \colorlet{grau}{sw-gray}
    \colorlet{hellgrau}{sw-white}
    \colorlet{loesung}{sw-dark}
    \colorlet{erfarbe}{sw-dark}
    \colorlet{irfarbe}{sw-dark}
\else
    \colorlet{spanisch}{spanisch-color}
    \colorlet{grau}{grau-color}
    \colorlet{hellgrau}{hellgrau-color}
    \colorlet{loesung}{loesung-color}
    \colorlet{erfarbe}{er-color}
    \colorlet{irfarbe}{ir-color}
\fi

\newcommand{\fachfarbe}{spanisch}

\newcommand{\thema}{Wiederholungstest: Mi tiempo libre}
\newcommand{\fach}{Spanisch}
\newcommand{\klasse}{7}
\newcommand{\maxpunkte}{30}
\newcommand{\zeit}{20 Minuten}
\newcommand{\datum}{\today}

\pagestyle{fancy}
\fancyhf{}
\renewcommand{\headrulewidth}{0.4pt}
\renewcommand{\footrulewidth}{0.4pt}

\lhead{\fach{} -- Klasse \klasse}
\chead{\textcolor{loesung}{\textbf{MT-05d MUSTERLÖSUNG}}}
\rhead{\datum}

\lfoot{\zeit}
\cfoot{}
\rfoot{Seite \thepage\ von \pageref{LastPage}}

\newcommand{\punktebox}[1]{\hfill\fbox{\small \_/#1}}
\newcommand{\luecke}[1]{\rule{#1}{0.4pt}}
\newcommand{\lsg}[1]{\textcolor{loesung}{\textbf{#1}}}

\newtcolorbox{infobox}{
    colback=hellgrau,
    colframe=\fachfarbe,
    boxrule=1pt,
    arc=0pt,
    left=8pt, right=8pt, top=8pt, bottom=8pt
}

\newenvironment{aufgabe}[1]{%
    \needspace{5\baselineskip}%
    \nopagebreak[4]%
    \vspace{10pt}
    \noindent\textbf{\textcolor{\fachfarbe}{Aufgabe #1}}
    \nopagebreak[4]%
    \vspace{5pt}
    \nopagebreak[4]%
    \begin{list}{}{\leftmargin=0pt}
    \item[]
}{%
    \end{list}
}

\newcommand{\es}[1]{\textcolor{\fachfarbe}{\textbf{#1}}}
\newcommand{\er}[1]{\textcolor{erfarbe}{\textbf{#1}}}
\newcommand{\ir}[1]{\textcolor{irfarbe}{\textbf{#1}}}

\begin{document}

% ═══════════════════════════════════════════════════════════════
% KOPFBEREICH
% ═══════════════════════════════════════════════════════════════

\begin{center}
    {\Large\bfseries\textcolor{\fachfarbe}{\thema}}\\[0.3em]
    {\large\textcolor{loesung}{\textbf{--- MUSTERLÖSUNG ---}}}
\end{center}

\vspace{5pt}

\noindent
\begin{tabularx}{\textwidth}{|X|X|X|}
\hline
\textbf{Name:} \lsg{Musterschüler/in} &
\textbf{Klasse:} \klasse &
\textbf{Datum:} \datum \\
\hline
\end{tabularx}

\vspace{10pt}

\begin{infobox}
\textbf{Zeit:} \zeit{} | \textbf{Punkte:} \maxpunkte{} BE | \textbf{Hilfsmittel:} keine
\end{infobox}

\vspace{15pt}

% ═══════════════════════════════════════════════════════════════
% AUFGABE 1: Hobbys (6P)
% ═══════════════════════════════════════════════════════════════

\begin{aufgabe}{1 -- Hobbys}
Wähle das passende Verb: \es{jugar}, \es{tocar}, \es{hacer}, \es{escuchar}, \es{ver}, \es{leer}. \punktebox{6}
\end{aufgabe}

\noindent
a) \lsg{jugar} al fútbol\\[0.8em]
b) \lsg{tocar} la guitarra\\[0.8em]
c) \lsg{hacer} deporte\\[0.8em]
d) \lsg{escuchar} música\\[0.8em]
e) \lsg{ver} la tele\\[0.8em]
f) \lsg{leer} un libro

\textit{\textcolor{grau}{Je 1P. jugar a = Spiel, tocar = Instrument, hacer = Aktivität.}}

\vspace{15pt}

% ═══════════════════════════════════════════════════════════════
% AUFGABE 2: Verben konjugieren (8P)
% ═══════════════════════════════════════════════════════════════

\begin{aufgabe}{2 -- Verben konjugieren}
Setze die richtige Form ein. \punktebox{8}
\end{aufgabe}

\noindent
a) Yo \lsg{como} una pizza. \er{(comer)}\\[1em]
b) Tú \lsg{bebes} agua. \er{(beber)}\\[1em]
c) Nosotros \lsg{vivimos} en Rostock. \ir{(vivir)}\\[1em]
d) Ella \lsg{escribe} un mensaje. \ir{(escribir)}

\textit{\textcolor{grau}{Je 2P. -er: yo -o, tú -es. -ir: nosotros -imos (nicht -emos!), ella -e.}}

\vspace{15pt}

% ═══════════════════════════════════════════════════════════════
% AUFGABE 3: gustar (8P)
% ═══════════════════════════════════════════════════════════════

\begin{aufgabe}{3 -- gustar}
Ergänze mit Pronomen + gusta/gustan. \punktebox{8}
\end{aufgabe}

\noindent
a) A mí \lsg{me gusta} nadar. (mir gefällt)\\[1em]
b) A ti \lsg{te gustan} los videojuegos. (dir gefallen)\\[1em]
c) A María \lsg{le gusta} el fútbol. (ihr gefällt)\\[1em]
d) A nosotros \lsg{nos gustan} las películas. (uns gefallen)

\textit{\textcolor{grau}{Je 2P. 1P Pronomen (me/te/le/nos), 1P gusta/gustan. Plural = gustan!}}

\vspace{15pt}

% ═══════════════════════════════════════════════════════════════
% AUFGABE 4: Übersetzung (8P)
% ═══════════════════════════════════════════════════════════════

\begin{aufgabe}{4 -- Übersetzung}
Übersetze ins Spanische. \punktebox{8}
\end{aufgabe}

\noindent
a) Was machst du gerne?

\lsg{¿Qué te gusta hacer?}

\vspace{1em}

\noindent
b) Ich höre gerne Musik.

\lsg{Me gusta escuchar música.}

\vspace{1em}

\noindent
c) Wir leben in Deutschland.

\lsg{(Nosotros) vivimos en Alemania.}

\vspace{1em}

\noindent
d) Er isst eine Pizza.

\lsg{(Él) come una pizza.}

\textit{\textcolor{grau}{Je 2P. Subjektpronomen optional. Häufige Fehler: *Yo gusto... (falsch!)}}

% ═══════════════════════════════════════════════════════════════
% PUNKTEÜBERSICHT
% ═══════════════════════════════════════════════════════════════

\vspace{25pt}
\hrule
\vspace{10pt}

\noindent
\begin{tabularx}{\textwidth}{|X|c|c|c|c||c|}
\hline
\textbf{Aufgabe} & 1 & 2 & 3 & 4 & \textbf{Gesamt} \\
\hline
\textbf{max. Punkte} & 6 & 8 & 8 & 8 & \textbf{30} \\
\hline
\textbf{erreicht} & \lsg{6} & \lsg{8} & \lsg{8} & \lsg{8} & \lsg{30} \\[0.8em]
\hline
\end{tabularx}

\vspace{10pt}

\begin{flushright}
\textbf{Note:} \lsg{14 NP (100\%)}
\end{flushright}

\end{document}
